\chapter{Metrics Analysis}
\label{ch:metricsAnalysis}

This chapter defines the five metrics used in this study
and evaluates their effectiveness in detecting the \acp{SSCA} discussed in \cref{ch:AnalysisRecentAttacks}.
\cref{sec:metricsMethodology} describes the methodology, 
how each metric is analyzed, how the baseline is established, and how to read the figures.
\cref{sec:analysisFileTypes} through \cref{sec:analysisEthAdd}
each cover one metric in full, from definition to comparison with the Malware dataset.
\cref{sec:wrapUpAnalysisMetrics} synthesizes the results and discusses the overall effectiveness, 
limitations, and cost of the proposed approach.

\section{Methodology}
\label{sec:metricsMethodology}

In this context, a metric is a measure computed from a specific version of an \ac{NPM} package.

While the literature offers a broad range of metrics based on code analysis~\cite{zoratti2025valutazione, weissberg2025llm},
we focused on a set of metrics defined ad hoc derived from the most relevant characteristics of the four \acp{SSCA} selected as case studies.
Metrics unrelated to these attacks, such as the total number of files in the tarball,
or \ac{API}-call patterns hidden behind obfuscation, were excluded.
The five selected metrics are: 
\emph{File Types} (\cref{sec:analysisFileTypes}),
\emph{Package Size} (\cref{sec:analysisPackageSize}),
\emph{Longest Line Length} (\cref{sec:analysisLLL}),
\emph{Unique Hexadecimal Values} (\cref{sec:analysisHex}),
and \emph{Unique Ethereum Addresses} (\cref{sec:analysisEthAdd}).
Each metric is discussed in a dedicated section:
first, its definition and the motivation behind its choice are presented;
then two analyses follow, structured as follows.

\subsection{Analysis of the Evolution of a Package}
\label{subsec:analysisEvolutionPackage}

We analyzed the distribution and evolutionary patterns of the metric within the Goodware dataset (\cref{subsec:goodwareDataset}),
with the goal of establishing a baseline for normal development conditions.
Understanding how a package evolves across versions is fundamental to evaluating potential attacks.
Tracking a metric across consecutive releases allows identifying anomalies that would be invisible in a single-version snapshot. 
A sudden deviation from an otherwise stable trend is precisely the signature left by the \acp{SSCA} analyzed in \cref{ch:AnalysisRecentAttacks}.
Moreover, this longitudinal perspective helps distinguish malicious injections from legitimate development events, 
such as the addition of build artifacts or a major refactoring, 
that can also produce outlier values but are traceable to documented development decisions.

Except for the \emph{File Types} metric (\cref{sec:analysisFileTypes}), 
which requires specialized representations, all remaining metrics share the same figure format for the Goodware dataset analysis.
A concrete example of this figure is \cref{fig:fig1PackageSize},
which shows the evolution of the Package Size metric across 20 versions.
The x-axis shows the package versions from the least recent (-19) to the latest available (0), 
and the y-axis shows the value of the metric on a logarithmic scale.
The vertical black bars mark the 95th-percentile confidence intervals, computed following an established statistical approach~\cite{article2}:
packages within the interval exhibit values consistent with normal behavior, 
while the gray dots above the bar represent the outlier values of the remaining 5\% of the population.
The red dashed line is the median, calculated for each package version and used instead of the mean to avoid distortion from extreme values.
Packages of particular interest are highlighted with colored lines,
for instance, \cref{fig:fig1PackageSizeIncreasing} (discussed in \cref{sec:analysisPackageSize}) 
shows the \texttt{moment-timezone} package (solid green) exhibiting a sudden jump at point~2, caused by the accidental inclusion of a temporary folder;
while in \cref{fig:fig1PackageSizeStable} \texttt{next} (solid orange) is a stable outlier across all versions
due to the large build artifacts it bundles by design.
Comparing these two cases already illustrates the core challenge of this analysis:
both produce outlier values, but for entirely different reasons,
and distinguishing a legitimate anomaly from a malicious one requires examining
the nature and history of the change.

\subsection{Effectiveness Analysis Against the Malware Dataset}
\label{subsec:analysisEffectivenessMetric}

Finally, we evaluated the effectiveness of each metric by comparing the values of the
compromised versions in the Malware dataset (\cref{subsec:malwareDataset}) against the Goodware baseline established above.
A metric is considered effective for a given attack if the compromised version produces
a value that falls among the outlier values of the Goodware distribution,
making it distinguishable from a legitimate release.
If the malicious version falls within the confidence interval,
the metric is not indicative of that attack,
since the compromised version is indistinguishable from legitimate ones.

Except for the \emph{File Types} metric, all remaining metrics share the same figure format for this comparison.
Two concrete examples are \cref{fig:fig2PackageSizeAttackD} and \cref{fig:fig2PackageSizeAttackB}, both for the \emph{Package Size} metric.
In \cref{fig:fig2PackageSizeAttackD}, the red crosses, marking the compromised versions of packages from Attack~D, 
fall among the outlier dots of the Goodware distribution, showing that the metric successfully detects this attack.
In \cref{fig:fig2PackageSizeAttackB}, by contrast, the red crosses remain within the confidence intervals; 
the malicious versions of Attack~B are indistinguishable from legitimate ones under this metric, 
which is therefore not indicative of that attack.
In both figures, the colored lines show the evolution of the compromised packages across the versions, 
providing the longitudinal context needed to see how sudden the anomaly introduced by the malicious version actually is.
When packages from the same attack overlap visually, separate figures are provided
per package for clarity, as in \cref{fig:fig2HexAttackA}.